\documentclass[11pt]{article}

\usepackage[margin=1in]{geometry}
\usepackage{booktabs}
\usepackage{longtable}
\usepackage{array}
\usepackage{pgf}
\usepackage{siunitx}
\usepackage{xcolor}
\usepackage[hidelinks]{hyperref}

\setlength{\emergencystretch}{3em}

\sisetup{
  detect-all,
  round-mode=places,
  round-precision=2,
  group-minimum-digits=4
}

\title{Drone-Carried OAI DU Dimensioning Model}
\author{OAI CU/DU Split Lab}
\date{2026-07-06}

% ---------------------------------------------------------------------------
% Editable mission assumptions.
% Change these values first when resizing the airborne payload.
% ---------------------------------------------------------------------------
\pgfmathsetmacro{\MissionMinutes}{20}
\pgfmathsetmacro{\ElectricalReserveFactor}{1.30}
\pgfmathsetmacro{\UsableBatteryFraction}{0.80}
\pgfmathsetmacro{\DCEfficiency}{0.88}
\pgfmathsetmacro{\BatterySpecificEnergyWhKg}{190}
\pgfmathsetmacro{\DronePayloadUtilization}{0.70}

% ---------------------------------------------------------------------------
% Editable configuration profiles.
% Masses are grams. Powers are watts. Each profile includes a dedicated payload
% electronics battery; if the drone supplies payload power, set the profile
% battery calculation to zero and keep a DC conversion allowance.
% ---------------------------------------------------------------------------
\pgfmathsetmacro{\MassPiFive}{46}
\pgfmathsetmacro{\MassJetsonCarrier}{250}
\pgfmathsetmacro{\MassMiniPc}{650}
\pgfmathsetmacro{\MassBTwoZeroFiveMini}{24}
\pgfmathsetmacro{\MassBTwoTen}{350}
\pgfmathsetmacro{\MassQuectelKit}{180}
\pgfmathsetmacro{\MassRfTimingLight}{180}
\pgfmathsetmacro{\MassRfTimingBTwoTen}{250}
\pgfmathsetmacro{\MassPowerMountLight}{300}
\pgfmathsetmacro{\MassPowerMountMedium}{350}
\pgfmathsetmacro{\MassPowerMountHeavy}{500}
\pgfmathsetmacro{\MassDualRadioExtra}{650}

\pgfmathsetmacro{\PowerPiFive}{18}
\pgfmathsetmacro{\PowerJetson}{25}
\pgfmathsetmacro{\PowerMiniPc}{45}
\pgfmathsetmacro{\PowerBTwoZeroFiveMini}{4}
\pgfmathsetmacro{\PowerBTwoTen}{8}
\pgfmathsetmacro{\PowerQuectel}{9}
\pgfmathsetmacro{\PowerAuxLight}{6}
\pgfmathsetmacro{\PowerAuxMedium}{8}
\pgfmathsetmacro{\PowerAuxHeavy}{10}
\pgfmathsetmacro{\PowerDualRadioExtra}{28}

% ---------------------------------------------------------------------------
% Derived profile calculations.
% ---------------------------------------------------------------------------
\pgfmathsetmacro{\WhFactor}{\MissionMinutes/60*\ElectricalReserveFactor/(\DCEfficiency*\UsableBatteryFraction)}

\pgfmathsetmacro{\AStaticG}{\MassPiFive+\MassBTwoZeroFiveMini+\MassQuectelKit+\MassRfTimingLight+\MassPowerMountLight}
\pgfmathsetmacro{\APowerW}{\PowerPiFive+\PowerBTwoZeroFiveMini+\PowerQuectel+\PowerAuxLight}
\pgfmathsetmacro{\ABatteryWh}{\APowerW*\WhFactor}
\pgfmathsetmacro{\ABatteryG}{(\ABatteryWh/\BatterySpecificEnergyWhKg)*1000}
\pgfmathsetmacro{\ATotalG}{\AStaticG+\ABatteryG}
\pgfmathsetmacro{\ARequiredPayloadKg}{\ATotalG/1000/\DronePayloadUtilization}

\pgfmathsetmacro{\BStaticG}{\MassJetsonCarrier+\MassBTwoTen+\MassQuectelKit+\MassRfTimingBTwoTen+\MassPowerMountMedium}
\pgfmathsetmacro{\BPowerW}{\PowerJetson+\PowerBTwoTen+\PowerQuectel+\PowerAuxMedium}
\pgfmathsetmacro{\BBatteryWh}{\BPowerW*\WhFactor}
\pgfmathsetmacro{\BBatteryG}{(\BBatteryWh/\BatterySpecificEnergyWhKg)*1000}
\pgfmathsetmacro{\BTotalG}{\BStaticG+\BBatteryG}
\pgfmathsetmacro{\BRequiredPayloadKg}{\BTotalG/1000/\DronePayloadUtilization}

\pgfmathsetmacro{\CStaticG}{\MassMiniPc+\MassBTwoTen+\MassQuectelKit+\MassRfTimingBTwoTen+\MassPowerMountHeavy}
\pgfmathsetmacro{\CPowerW}{\PowerMiniPc+\PowerBTwoTen+\PowerQuectel+\PowerAuxHeavy}
\pgfmathsetmacro{\CBatteryWh}{\CPowerW*\WhFactor}
\pgfmathsetmacro{\CBatteryG}{(\CBatteryWh/\BatterySpecificEnergyWhKg)*1000}
\pgfmathsetmacro{\CTotalG}{\CStaticG+\CBatteryG}
\pgfmathsetmacro{\CRequiredPayloadKg}{\CTotalG/1000/\DronePayloadUtilization}

\pgfmathsetmacro{\DStaticG}{\CStaticG+\MassDualRadioExtra}
\pgfmathsetmacro{\DPowerW}{\CPowerW+\PowerDualRadioExtra}
\pgfmathsetmacro{\DBatteryWh}{\DPowerW*\WhFactor}
\pgfmathsetmacro{\DBatteryG}{(\DBatteryWh/\BatterySpecificEnergyWhKg)*1000}
\pgfmathsetmacro{\DTotalG}{\DStaticG+\DBatteryG}
\pgfmathsetmacro{\DRequiredPayloadKg}{\DTotalG/1000/\DronePayloadUtilization}

\newcommand{\numkg}[1]{\pgfmathparse{#1}\num{\pgfmathresult}}
\newcommand{\numg}[1]{\pgfmathparse{#1}\num{\pgfmathresult}}
\newcommand{\numw}[1]{\pgfmathparse{#1}\num{\pgfmathresult}}
\newcommand{\numwh}[1]{\pgfmathparse{#1}\num{\pgfmathresult}}
\newcommand{\fit}{\textcolor{green!45!black}{Fit}}
\newcommand{\tight}{\textcolor{orange!70!black}{Tight}}
\newcommand{\nofit}{\textcolor{red!70!black}{No}}

\begin{document}
\maketitle

\section*{Purpose}

This document turns the portable-DU idea into an editable sizing model. It does
not claim that the airborne deployment has been validated. The lab rollback
baseline remains Ethernet CU/DU with SIB8 at OAI commit
\texttt{102965a669b9444857c27843ec8ce62780bf9d37}; the drone work should be
treated as a future packaging and flight-integration step after the Quectel F1
path is stable.

Public component and drone specifications cited here were checked on
2026-07-06. Treat the cited aircraft ratings as planning inputs only; confirm
the exact regional manual, payload mount, firmware limitations, and supplier
configuration before procurement or flight.

The intended architecture is one ground OAI 5GC plus one shared CU, one donor DU
for the Quectel modem-serving cell, and one airborne access DU carrying the
local USRP access radio. The Quectel modem provides the wireless F1 backhaul to
the ground CU. The USRP B210 is the current lab access-radio baseline; the USRP
B205mini-i is listed as a lighter future radio candidate, not as a proven
drop-in replacement for the current B210 profile.

\section*{Editable Formula Model}

The sizing model uses component switches, mass assumptions, and power
assumptions. The equations below dimension the payload electronics pack and the
aircraft payload-rating margin. They are not a full rotorcraft endurance model:
the aircraft flight battery, rotor/motor efficiency, propeller geometry, air
density, wind, and actual flight path must still be checked with the selected
drone's vendor tools or flight tests.

\subsection*{Formula provenance}

The payload-electronics energy equation is a constant-power battery sizing
model. It follows the standard energy relation $E=P\,t$, with derating for
usable capacity and DC conversion efficiency. This is consistent with UAV
endurance literature that estimates endurance from available energy divided by
required power, then refines that estimate with battery-discharge behaviour.
Hwang, Cha, and Jung describe multirotor endurance estimation as a required
power calculation followed by battery-discharge calculation, and note that
endurance is commonly estimated from available energy and power consumption;
their practical method also includes payload weight in total vehicle
weight.\footnote{M.-h. Hwang, H.-R. Cha, and S. Y. Jung, ``Practical Endurance
Estimation for Minimizing Energy Consumption of Multirotor Unmanned Aerial
Vehicles,'' \emph{Energies}, vol. 11, no. 9, 2221, 2018.
\url{https://www.mdpi.com/1996-1073/11/9/2221}} Abdilla, Richards, and
Burrow derive battery-powered rotorcraft endurance formulae with available and
usable battery-capacity factors, which motivates keeping an explicit usable
battery fraction rather than using nominal pack energy directly.\footnote{A.
Abdilla, A. Richards, and S. Burrow, ``Power and Endurance Modelling of
Battery-Powered Rotorcraft,'' \emph{IEEE/RSJ International Conference on
Intelligent Robots and Systems workshop proceedings}, 2015.
\url{https://www.semanticscholar.org/paper/Power-and-endurance-modelling-of-battery-powered-Abdilla-Richards/76ffb97d686447bca76d36a079762e88c6f4bc30}}

The selected-drone gate
$C_\mathrm{required}=m_\mathrm{payload}/\gamma_\mathrm{payload}$ is not taken
from a paper. It is a conservative lab margin rule: the OAI payload should use
only a chosen fraction of the advertised drone payload rating. The literature
does justify why this margin is needed, because multirotor endurance and
required power are strongly affected by payload weight, drag, battery discharge,
and the selected flight regime.\footnote{Y. Zeng, J. Xu, and R. Zhang,
``Energy Minimization for Wireless Communication with Rotary-Wing UAV,''
\emph{IEEE Transactions on Wireless Communications}, vol. 18, no. 4, pp.
2329--2345, 2019. The paper separates communication-related energy from
propulsion energy and gives a rotary-wing propulsion power model.
\url{https://arxiv.org/pdf/1804.02238}}

For a selected configuration $S$:

\[
  m_\mathrm{fixed}(S)=\sum_i x_i(S)m_i
\]

\[
  P_\mathrm{payload}(S)=\sum_i x_i(S)P_i
\]

\[
  E_\mathrm{nom}(S)=
  \frac{P_\mathrm{payload}(S)t_\mathrm{mission}r_\mathrm{reserve}}
       {\eta_\mathrm{dc}u_\mathrm{battery}}
\]

\[
  m_\mathrm{battery}(S)=
  \frac{1000E_\mathrm{nom}(S)}{\rho_\mathrm{battery}}
\]

\[
  m_\mathrm{payload}(S)=m_\mathrm{fixed}(S)+m_\mathrm{battery}(S)
\]

\[
  C_\mathrm{required}(S)=
  \frac{m_\mathrm{payload}(S)}{\gamma_\mathrm{payload}}
\]

A drone is considered eligible when
$C_\mathrm{drone}\geq C_\mathrm{required}$. The default
$\gamma_\mathrm{payload}=0.70$ keeps the electronics payload below
\SI{70}{\percent} of the published drone payload rating, leaving margin for hot
weather, altitude, vibration isolation, mounting mistakes, antenna separation,
and small later additions.

\begin{table}[h]
\centering
\caption{Global sizing assumptions defined at the top of this \LaTeX{} file.}
\begin{tabular}{lll}
\toprule
Symbol & Current value & Meaning \\
\midrule
$t_\mathrm{mission}$ & \num{\MissionMinutes} min & Payload electronics run time \\
$r_\mathrm{reserve}$ & \num{\ElectricalReserveFactor} & Electrical reserve multiplier \\
$u_\mathrm{battery}$ & \num{\UsableBatteryFraction} & Usable fraction of nominal pack energy \\
$\eta_\mathrm{dc}$ & \num{\DCEfficiency} & DC/DC conversion efficiency \\
$\rho_\mathrm{battery}$ & \num{\BatterySpecificEnergyWhKg} Wh/kg & Practical LiPo/Li-ion pack specific energy \\
$\gamma_\mathrm{payload}$ & \num{\DronePayloadUtilization} & Max allowed fraction of drone payload rating \\
\bottomrule
\end{tabular}
\end{table}

\section*{Component Defaults}

The values below are deliberately conservative for early design. Vendor-published
board weights are used where available. Carrier boards, RF cabling, thermal
management, antennas, payload plate, fasteners, and power conversion are grouped
as editable engineering allowances because they depend on the mechanical design.

\begin{longtable}{p{0.25\linewidth}p{0.15\linewidth}p{0.13\linewidth}p{0.38\linewidth}}
\caption{Editable component assumptions.}\\
\toprule
Component & Mass & Power & Source and sizing note \\
\midrule
\endfirsthead
\toprule
Component & Mass & Power & Source and sizing note \\
\midrule
\endhead
Raspberry Pi 5 board & \numg{\MassPiFive} g & \numw{\PowerPiFive} W &
Board mass should be re-weighed with cooler and storage. Raspberry Pi lists the Pi 5 platform and the official supply is commonly sized around the 5 V, 5 A class.\footnote{\url{https://www.raspberrypi.com/products/raspberry-pi-5/}} \\
Jetson Orin Nano carrier class & \numg{\MassJetsonCarrier} g & \numw{\PowerJetson} W &
Editable mass for module, carrier, cooling, and storage. NVIDIA lists Jetson Orin Nano Super Developer Kit power modes from 7 W to 25 W.\footnote{\url{https://www.nvidia.com/en-us/autonomous-machines/embedded-systems/jetson-orin/nano-super-developer-kit/}} \\
Mini-PC/x86 class DU & \numg{\MassMiniPc} g & \numw{\PowerMiniPc} W &
Editable stand-in for the smallest x86 host capable of the chosen OAI real-time profile. ASUS NUC 13 Pro configurations list cTDP classes from 20 W to 45 W depending on CPU.\footnote{\url{https://www.asus.com/us/displays-desktops/nucs/nuc-mini-pcs/asus-nuc-13-pro/techspec/}} \\
USRP B205mini-i & \numg{\MassBTwoZeroFiveMini} g & \numw{\PowerBTwoZeroFiveMini} W &
Light SDR candidate. NI/Ettus B205mini-i published dimensions are 83.3 x 50.8 x 8.4 mm and weight is 24.0 g; it is bus-powered over USB 3.0.\footnote{\url{https://mm.digikey.com/Volume0/opasdata/d220001/medias/docus/822/6002-410-021_Web.pdf}} \\
USRP B210 & \numg{\MassBTwoTen} g & \numw{\PowerBTwoTen} W &
Current lab access-radio baseline. Ettus/Digilent lists 97 x 155 x 15 mm and 350 g for the B210.\footnote{\url{https://digilent.com/shop/ni-ettus-usrp-b210-2x2-70mhz-6ghz-sdr-cognitive-radio/}} \\
Quectel RM500Q-GL kit & \numg{\MassQuectelKit} g & \numw{\PowerQuectel} W &
Module plus M.2 carrier, thermal pad/heatsink, SIM holder, USB/PCIe adapter, and antennas. The module itself is about 8.7 g and 30.0 x 52.0 x 2.3 mm.\footnote{\url{https://www.icpdas.com/web/product/download/wireless/modem/3g4g/rm500q-gl/document/data_sheet/RM500Q-GL_datasheet_en.pdf}} \\
Light RF/timing set & \numg{\MassRfTimingLight} g & \numw{\PowerAuxLight} W &
Editable allowance for GNSS/timing, filters, short RF cables, antenna mounts, and small fan loads. \\
B210 RF/timing set & \numg{\MassRfTimingBTwoTen} g & \numw{\PowerAuxMedium} W &
Editable allowance for heavier antenna separation, coax, filters, clocking, and thermal margin. \\
Power/mount light & \numg{\MassPowerMountLight} g & included & DC regulators, payload plate, isolation, fasteners, enclosure, wiring, and safety tether. \\
Power/mount medium & \numg{\MassPowerMountMedium} g & included & Same allowance with stronger B210 vibration isolation. \\
Power/mount heavy & \numg{\MassPowerMountHeavy} g & included & Same allowance for x86 host, larger heatsink, and serviceable field enclosure. \\
\bottomrule
\end{longtable}

\section*{Computed Payload Configurations}

\begin{longtable}{p{0.18\linewidth}p{0.27\linewidth}rrrrr}
\caption{Computed electronics payload for the default \num{\MissionMinutes} minute mission.}\\
\toprule
ID & Configuration & Fixed mass g & Power W & Battery Wh & Battery g & Payload g \\
\midrule
\endfirsthead
\toprule
ID & Configuration & Fixed mass g & Power W & Battery Wh & Battery g & Payload g \\
\midrule
\endhead
A & Minimum airborne access-DU proof: Pi 5, USRP B205mini-i, Quectel kit & \numg{\AStaticG} & \numw{\APowerW} & \numwh{\ABatteryWh} & \numg{\ABatteryG} & \numg{\ATotalG} \\
B & B210 validation payload: Jetson-class compute, USRP B210, Quectel kit & \numg{\BStaticG} & \numw{\BPowerW} & \numwh{\BBatteryWh} & \numg{\BBatteryG} & \numg{\BTotalG} \\
C & Actual mini-PC class: x86 mini-PC, USRP B210, Quectel kit & \numg{\CStaticG} & \numw{\CPowerW} & \numwh{\CBatteryWh} & \numg{\CBatteryG} & \numg{\CTotalG} \\
D & Dual-radio prototype: mini-PC class plus extra donor/access RF margin & \numg{\DStaticG} & \numw{\DPowerW} & \numwh{\DBatteryWh} & \numg{\DBatteryG} & \numg{\DTotalG} \\
\bottomrule
\end{longtable}

\begin{table}[h]
\centering
\caption{Drone payload rating required when using the default 70 percent utilization rule.}
\begin{tabular}{lrr}
\toprule
Configuration & Computed payload kg & Required drone payload rating kg \\
\midrule
A & \numkg{\ATotalG/1000} & \numkg{\ARequiredPayloadKg} \\
B & \numkg{\BTotalG/1000} & \numkg{\BRequiredPayloadKg} \\
C & \numkg{\CTotalG/1000} & \numkg{\CRequiredPayloadKg} \\
D & \numkg{\DTotalG/1000} & \numkg{\DRequiredPayloadKg} \\
\bottomrule
\end{tabular}
\end{table}

\section*{Drone Selection Matrix}

Payload ratings change by battery choice, temperature, altitude, propellers, and
regulatory takeoff weight limits. The values below are planning gates, not
purchase approval. Re-check the current manual for the exact region and payload
mount before buying or flying.

\begin{longtable}{p{0.20\linewidth}p{0.16\linewidth}p{0.18\linewidth}p{0.27\linewidth}}
\caption{Candidate drones for each configuration.}\\
\toprule
Drone & Published payload gate & Fits default profiles & Recommendation \\
\midrule
\endfirsthead
\toprule
Drone & Published payload gate & Fits default profiles & Recommendation \\
\midrule
\endhead
Freefly Astro Max & 3.0 kg maximum payload; Freefly also lists about 18 min at 3.0 kg payload.\footnote{\url{https://freeflysystems.com/astro/specs}}\footnote{\url{https://freeflysystems.com/knowledge-base/how-long-can-astro-fly}} &
A: \fit, B: \fit, C: \nofit, D: \nofit &
Best compact research drone if the payload is kept near the light or Jetson/B210 profile. Not the first choice for x86 mini-PC packaging. \\
DJI Matrice 350 RTK & About 2.7 kg gross-weight margin from 9.2 kg max takeoff weight minus 6.47 kg with TB65 batteries; the single downward gimbal damper is rated much lower at 960 g.\footnote{\url{https://enterprise.dji.com/matrice-350-rtk/specs}} &
A: \fit, B: \fit, C: \nofit, D: \nofit &
Good mature enterprise platform for the light and B210 validation payloads only if the custom mount carries the load through an approved airframe interface rather than the small single-gimbal damper. \\
DJI Matrice 400 & Up to 6 kg payload and 59 min maximum flight time.\footnote{\url{https://enterprise.dji.com/matrice-400/specs}} &
A: \fit, B: \fit, C: \fit, D: \fit &
Best first serious choice for the actual x86 mini-PC plus B210 plus Quectel configuration. It keeps enough payload margin for a real enclosure and later RF changes. \\
Inspired Flight IF1200A & 8.6 kg max payload class, with 24.9 kg max gross takeoff weight in current specifications.\footnote{\url{https://www.inspiredflight.com/if1200}} &
A: \fit, B: \fit, C: \fit, D: \fit &
Heavy-lift research choice when the payload needs open integration space, larger batteries, or test instrumentation beyond the radio stack. \\
Skyfront Perimeter 8 / 8+ hybrid & Skyfront advertises 5 kg payload for 3 h and 10 kg maximum payload for 1 h on the Perimeter 8 family.\footnote{\url{https://skyfront.com/}} &
A: \fit, B: \fit, C: \fit, D: \fit &
Use only when endurance is the primary requirement. It is mechanically and operationally larger than needed for first airborne DU proofs. \\
\bottomrule
\end{longtable}

\section*{Configuration Decisions}

\subsection*{Recommended first build}

Start with Configuration B if the priority is continuity with the current lab
radio baseline. It keeps the USRP B210, uses a smaller compute class than the
current mini-PC, and remains eligible for both DJI Matrice 350 RTK and DJI
Matrice 400 with the default margin rule. The Matrice 400 is the cleaner first
purchase target because it leaves meaningful reserve for enclosure growth,
antenna standoffs, RF shielding, heat, and a second modem or timing board.

\subsection*{Minimum-weight research path}

Configuration A is the best mechanical path if the USRP B205mini-i proves
sufficient for the selected OAI bandwidth and access-cell objective. It reduces
the radio board mass from the B210's 350 g class to about 24 g, but it should be
treated as an experiment until OAI attach, PWS/SIB8 observation, user-plane
traffic, Quectel F1 traffic, and Ethernet rollback are all recorded.

\subsection*{Actual mini-PC path}

Configuration C is the closest to carrying the current mini-PC style deployment
without pretending the packaging is already optimized. It should skip the
Matrice 350/Astro class and go directly to a 6 kg-class or heavier drone. For
this configuration, the Matrice 400 is the most balanced option; IF1200A is the
reserve-heavy alternative.

\subsection*{Dual-radio or full donor/access prototype}

Configuration D models the case where the airborne payload grows beyond one
access radio and one modem, for example by adding extra RF, a second SDR path,
or field instrumentation. It belongs on Matrice 400 or heavier platforms from
the beginning. The IF1200A or Skyfront class becomes attractive only if payload
development, endurance, or integration access matters more than portability.

\section*{Pre-Flight Engineering Checklist}

\begin{itemize}
  \item Weigh the exact compute board, SDR, Quectel carrier, antennas, coax,
        heatsinks, payload battery, DC converters, enclosure, and mount before
        selecting the final drone.
  \item Recompute this document with the measured masses and with the real OAI
        power draw under radio load, not idle desktop power.
  \item Keep the USRP B210 as the verified access-radio baseline until the
        lighter USRP-mini path has fresh attach, PWS, user-plane, and rollback
        evidence.
  \item Keep generated configs, modem logs, packet captures, SIM data, WireGuard
        keys, and subscriber material out of Git.
  \item Record sanitized evidence for every airborne trial: OAI commit, drone
        payload mass, power profile, thermal state, F1-C/F1-U packet proof,
        Quectel state, phone service/PWS observation, and rollback test.
  \item Re-check local flight approval, spectrum authorization, payload mounting
        limits, propeller clearance, and battery transport rules before any
        outdoor flight.
\end{itemize}

\section*{How to Re-Use This File}

Edit the numeric variables near the top of this file, then compile:

\begin{verbatim}
tectonic docs/drone-du-dimensioning.tex
\end{verbatim}

The tables will update automatically. The most important values to change after
the first bench packaging pass are \texttt{\textbackslash MissionMinutes},
\texttt{\textbackslash BatterySpecificEnergyWhKg},
\texttt{\textbackslash DronePayloadUtilization}, component masses, and component
powers.

\end{document}
